\documentclass[exercice,a4]{thedocument}%
\usepackage{texmaths-tikz}%
\usepackage{tabularx}%
\usepackage{tabularray}%
\usepackage[np]{numprint}%
\startdatakeys
\source{24GENMATG11}
\BO{2026}
\cycle{4}
\competencesDuSocle{CA,CO}
\objectifsApprentissage{C4A3bc,C4D2aa,C4C1bc,C4C2ab,C4C3aa}
\motsCles{écriture scientifique,vitesse,probabilité,moyenne,médiane,fraction}
\enddatakeys
\begin{document}%
Cet exercice est un questionnaire à choix multiples (QCM). Aucune
justification n'est demandée. Pour chaque question, quatre réponses sont
proposées : une seule réponse est exacte.

\begin{enumerate}
    \item Donner l'écriture scientifique de $0,193\times10^{-100}$.

    \begin{tblr}{hlines,vlines,colspec={*{4}{X[c]}}}
        $1,93\times10^{-99}$ & $1,93\times10^{-101}$ &
        $193\times10^{-103}$ & $193\times10^{-97}$
    \end{tblr}

    \item Lili parcourt $\np[km]{480}$ en $\np[h]{5}~\np[min]{42}$.
    Quelle est sa vitesse moyenne en $\mathrm{km}/\mathrm{h}$, arrondie au
    dixième ?

    \begin{tabularx}{\linewidth}{|*{4}{>{\centering\arraybackslash}X|}}\hline
        88,6 & 84,2 & 1,4 & 23,4\\ \hline
    \end{tabularx}

    \item Une roue comporte quinze secteurs équiprobables. Les nombres
    indiqués sont, dans l'ordre, $2$, un nombre effacé, $2$, $1$, $1$, $2$,
    $2$, $1$, $2$, $1$, $1$, $1$, $2$, $2$, $2$.
    Est-il possible d'écrire un nombre dans le secteur effacé pour que la
    probabilité d'obtenir 2 soit $\dfrac35$ ?

    \begin{tabularx}{\linewidth}{|*{4}{>{\centering\arraybackslash}X|}}\hline
        Oui, en écrivant 1 & Oui, en écrivant 2 &
        Ce n'est pas possible & Oui, en laissant le secteur vide\\ \hline
    \end{tabularx}

    \item Pour la liste $5~;~1~;~3~;~10~;~17~;~11~;~10$, que représente le
    nombre 5 ?

    \begin{tabularx}{\linewidth}{|*{4}{>{\centering\arraybackslash}X|}}\hline
        La médiane & L'étendue & La moyenne & Rien de particulier\\ \hline
    \end{tabularx}

    \item Léa paye $\dfrac15$ du prix de son vélo, puis le reste en trois
    paiements égaux. Quelle fraction du prix représente l'un de ces paiements ?

    \begin{tblr}{hlines,vlines,colspec={*{4}{X[c]}}}
        $\dfrac{12}{5}$ & $\dfrac1{15}$ & $\dfrac4{15}$ & $\dfrac35$
    \end{tblr}
\end{enumerate}
\showthesource
\end{document}
